\documentclass{article}
\usepackage[utf8]{inputenc}
\usepackage{graphicx}
\usepackage{amsmath}
\usepackage{hyperref}
\usepackage{authblk}
\usepackage{geometry}
\geometry{margin=1in}

\title{Basket of Fruit Recognition Using Deep Learning}
\author[1]{Heran E}
\author[1]{Iman I}
\author[1]{Ruhama Y}
\author[1]{Samrawit K}
\author[1]{Yordanos M}
\affil[1]{School of Information Technology and Engineering, Addis Ababa University, Ethiopia}
\date{January 2026}

\begin{document}

\maketitle

\begin{abstract}
Automated fruit recognition is an important task in computer vision with applications in agriculture (sorting and harvesting) and retail (smart checkout) systems. We present a custom 5-block convolutional neural network (CNN) trained \textbf{from scratch} to classify 10 fruit types. To handle real-world variability, we collected a dataset of several thousand images captured under diverse lighting, background, and angle conditions. The data pipeline includes standard preprocessing and extensive augmentation: random flips, rotations, zooms, translations, and MixUp. Training uses the Adam optimizer, a cosine-decayed learning rate, and strong regularization techniques. The model achieves a \textbf{peak validation accuracy of 85\%} and \textbf{test accuracy of 80.17\%}. These results show that a moderately deep CNN can effectively classify fruits without transfer learning, providing a strong baseline for deployment on resource-constrained systems.
\end{abstract}

\section{Introduction}
Convolutional Neural Networks (CNNs) have become the dominant tool for image classification tasks in computer vision. Their ability to extract hierarchical representations directly from raw image data has significantly reduced the need for manual feature engineering. While models such as ResNet or MobileNet provide excellent performance due to pretraining on large-scale datasets like ImageNet, they are often too resource-intensive for real-world applications, especially in embedded or mobile settings where compute and memory are limited.

Fruit recognition serves as a practical testbed for evaluating CNN architectures because of its relevance to agricultural automation and retail systems. Despite seemingly distinct categories, this problem presents real challenges due to intra-class variability (e.g., ripeness and size) and inter-class similarity (e.g., lemons and oranges under similar lighting). To address this, our project explores the training of a custom CNN from scratch on a small, custom-collected dataset of fruit images.

This work demonstrates that with carefully chosen data augmentation strategies and regularization techniques, a moderate-capacity CNN can achieve competitive accuracy. In this paper, we detail the full pipeline from dataset collection to evaluation and offer insights into the limitations and strengths of training from scratch in a constrained setting.

\section{Related Work}
Early attempts at fruit classification relied on hand-crafted features such as color histograms, texture descriptors, and contour-based shape analysis, often in combination with traditional machine learning classifiers like SVMs or k-NN. These approaches, while effective in controlled environments, struggle in diverse real-world conditions with complex backgrounds and lighting variability.

The advent of deep learning, particularly with AlexNet~\cite{krizhevsky2012imagenet}, revolutionized image classification by introducing CNNs trained end-to-end on large datasets. In fruit recognition, many researchers have since adopted transfer learning, fine-tuning pretrained models such as VGG16, ResNet50, or MobileNet on fruit datasets. These methods significantly reduce training time and typically yield high accuracy, even on limited data.

However, such approaches come with dependencies on large external datasets and pretraining infrastructure. In contrast, our approach aligns with research focusing on lightweight CNNs trained from scratch, particularly in low-resource settings. Additionally, Zhang et al.~\cite{zhang2017mixup} introduced MixUp, a regularization technique that interpolates between random sample pairs during training, which has been shown to improve model generalization. Our study builds on this work by applying MixUp in combination with geometric augmentations in a from-scratch training setup.

\section{Problem Definition and Dataset}
\subsection{Problem Statement}
This work aims to develop a robust classifier capable of categorizing images of fruits into one of ten predefined classes: Apple, Avocado, Banana, Lemon, Mango, Orange, Papaya, Pineapple, Tomato, and Watermelon. The classifier operates on static RGB images, each containing a single fruit captured in a natural setting. The classification system is intended to function reliably across variations in lighting, camera angle, and fruit orientation.

\subsection{Dataset Collection and Split}
Our dataset was collected manually using smartphone cameras under diverse lighting and background conditions to emulate real-world scenarios such as homes and local markets. This diversity introduces natural challenges like shadows, occlusion, and clutter, making the classification task more realistic.

The dataset comprises approximately 4000 labeled images, evenly distributed across the ten fruit classes. We split the dataset into three subsets: 70\% for training to fit model parameters, 15\% for validation to monitor performance during training, and 15\% for final testing to assess generalization. Each subset maintains class balance to ensure fair evaluation, as shown in Figure \ref{fig:distribution}.

\begin{figure}[htbp]
    \centering
    \includegraphics[width=0.8\textwidth]{results/class_distribution.png}
    \caption{Distribution of fruit classes in the training dataset.}
    \label{fig:distribution}
\end{figure}

\subsection{Preprocessing and Augmentation}
All images are resized to 192x192 pixels to standardize input dimensions. Pixel intensities are scaled to the [0, 1] range and normalized using mean and standard deviation values computed from the training set. To improve model robustness, we apply several augmentation strategies during training:

We use random horizontal flips to simulate mirrored viewpoints, small-angle rotations (up to ±3°) to account for camera angle variation, and slight zoom and translation to mimic different distances and positions. Most importantly, we employ MixUp with a mixing parameter $\alpha=0.2$, which generates synthetic samples by linearly combining image-label pairs. This technique helps the model learn smoother decision boundaries and reduces overfitting.

\section{Model and Training Methodology}
\subsection{CNN Architecture}
The CNN architecture used in this study consists of five sequential convolutional blocks designed to progressively extract higher-level features from the input images. Each of the first four blocks includes two convolutional layers followed by batch normalization and ReLU activation, with a max pooling layer to downsample spatial dimensions.

Specifically, Block 1 uses 32 filters, Block 2 uses 64, Block 3 uses 128, and Block 4 uses 256. The fifth block contains a single convolutional layer with 320 filters, followed by batch normalization and pooling. After the convolutional backbone, we apply Global Average Pooling to reduce the feature maps to a single vector, followed by Dropout (0.45) to mitigate overfitting. The final classifier includes a dense layer with 384 units and L2 regularization, culminating in a 10-unit softmax layer for multiclass output. The complete architecture is visualized in Figure \ref{fig:architecture}.

\begin{figure}[htbp]
    \centering
    \includegraphics[width=0.6\textwidth]{results/cnn_architecture.png}
    \caption{Proposed Stronger\_CNN architecture with 5 convolutional blocks.}
    \label{fig:architecture}
\end{figure}

\subsection{Training Details}
The model is trained using the Adam optimizer, which combines the benefits of momentum and adaptive learning rates. We set the initial learning rate to $2\times10^{-4}$ and employ a cosine decay schedule to gradually reduce the learning rate over time, promoting stable convergence. Training is performed for up to 70 epochs with a batch size of 32. We apply early stopping with a patience of 10 epochs, monitoring validation accuracy to avoid overfitting.

The categorical crossentropy loss function is used to measure the discrepancy between predicted and true class distributions. Regularization strategies include dropout in the classifier head (rate 0.45), L2 weight decay on the dense layer ($2\times10^{-4}$), and MixUp augmentation during training.

\subsection{Implementation Environment}
We implemented our training pipeline using TensorFlow 2.x and Keras. All experiments were conducted on an NVIDIA GPU for efficient computation. Data loading and preprocessing were performed using TensorFlow's `tf.data` API, which allows real-time augmentation and efficient batch preparation.

\section{Results and Discussion}
\subsection{Evaluation Metrics}
To evaluate our model, we report accuracy and crossentropy loss on both validation and test sets. The best performing model, as selected by highest validation accuracy, achieved 85.00\% accuracy on the validation set and 80.17\% on the unseen test set, with a final test loss of approximately 1.04. The training progression is shown in Figure \ref{fig:curves}.

\begin{figure}[htbp]
    \centering
    \includegraphics[width=1.0\textwidth]{results/accuracy_loss_curves.png}
    \caption{Training and validation accuracy and loss over 70 epochs.}
    \label{fig:curves}
\end{figure}

\subsection{Observations}
These results indicate that the model has learned robust features from the training data and generalizes reasonably well to new samples. While accuracy is strong, certain class-level confusions persist. For instance, the model occasionally misclassifies green apples as avocados or lemons as oranges, likely due to color and shape similarity. Lighting artifacts and background clutter also affect predictions. The relatively small gap between validation and test performance suggests the model is not overfitting excessively. Detailed performance is shown in the confusion matrix (Figure \ref{fig:cm}) and qualitative examples are provided in Figure \ref{fig:examples}.

\begin{figure}[htbp]
    \centering
    \includegraphics[width=0.8\textwidth]{results/confusion_matrix.png}
    \caption{Confusion matrix highlighting class-level performance on the test set.}
    \label{fig:cm}
\end{figure}

\begin{figure}[htbp]
    \centering
    \includegraphics[width=1.0\textwidth]{results/prediction_examples.png}
    \caption{Sample predictions from the test set with true vs. predicted labels.}
    \label{fig:examples}
\end{figure}

Overall, our architecture combined with data augmentation and regularization strategies provides a solid baseline. The results demonstrate that even in the absence of transfer learning, effective fruit classification is possible with thoughtful design and tuning.

\section{Conclusion and Future Work}
In this paper, we presented a complete deep learning pipeline for fruit image classification using a CNN trained from scratch. Our system achieves a test accuracy of over 80\%, confirming that moderate CNNs can perform well on small, diverse datasets. We have shown that techniques like MixUp, dropout, and learning rate scheduling significantly enhance performance.

For future work, we plan to compare our from-scratch model with pretrained networks such as EfficientNet and MobileNet to establish upper-bound accuracy. We also aim to compress and quantize the model for edge deployment, expand the dataset with more fruit classes, and explore transformer-based architectures that have recently shown strong performance in image classification.

\bibliographystyle{plain}
\bibliography{references}

\end{document}
